\subsection{测度的定义}

定义 2. --- \(\mathcal{K}(\mathrm{X};\mathbf{C})\) （\(\mathrm{X}\) 为一个局部紧空间）上的一个连续线性形式称为 \(\mathrm{X}\) 上的一个测度（或复测度）。

若 \(\mu\) 是局部紧空间 \(X\) 上的一个测度，则测度对一个函数 \(f \in \mathcal{K}(X; \mathbf{C})\) 所取的值称为 \(f\) 关于 \(\mu\) 的积分；除了通用记号 \(\mu(f)\) 与 \(\langle f, \mu \rangle\) 之外，也使用记号 \(\int f d\mu, \int f\mu, \int f(x) d\mu(x)\) 与 \(\int f(x)\mu(x)\) 来表示它；关于字母 \(x\) 的用法，参见 S, I, §1, No. 1。

根据直接极限中连续性的判别法（TVS, Ch. II, §4, No. 4, 命题 5），说 \(\mu\) 是 X 上的一个测度，就是说 \(\mu\) 是 \(\mathcal{K}(\mathrm{X};\mathbf{C})\) 上满足下述条件的线性形式：对 X 的每个紧子集 K，存在一个数 \(\mathrm{M_K}\) ，使得对每个支集含于 K 的函数 \(f\in \mathcal{K}(\mathrm{X};\mathbf{C})\) ，

\[
| \mu (f) | \leqslant \mathrm{M} _ {\mathrm{K}} \cdot \| f \| \quad (\text { where } \| f \| = \sup _ {x \in \mathrm{X}} | f (x) |).\tag{4}
\]

更一般地：

命题 6. --- 设 X 是一个局部紧空间，\((\mathrm{K}_{\alpha})\) 是 X 的一族紧子集，它们的内部 \(\mathring{K}_{\alpha}\) 构成 X 的一个开覆盖。为使线性形式 \(\mu\) （\(\mathcal{K}(\mathrm{X};\mathbf{C})\) 上的）是 X 上的测度，必须且只须对每个 \(\alpha\) 存在一个数 \(M_{\alpha}\) ，使得

\[
| \mu (f) | \leqslant \mathrm{M} _ {\alpha} \cdot \| f \|\tag{5}
\]

对每个函数 \(f \in \mathcal{K}(\mathrm{X}, \mathrm{K}_{\alpha}; \mathbf{C})\) 成立。

这一条件显然是必要的，故只须证明 (5) 对 X 的每个紧子集 K 蕴涵 (4)。现在，K 被有限个开集 \(\mathring{K}_{\alpha_{i}}\) （ \(1 \leqslant i \leqslant n\) ）所覆盖；把 No. 2 的引理 1 应用于 K 与诸 \(\mathring{K}_{\alpha_{i}}\) ，存在 X 上的连续函数 \(g_{i} \geqslant 0\) ，使得 \(\operatorname{Supp}(g_{i}) \subset \mathrm{K}_{\alpha_{i}}\) ，\(0 \leqslant \sum_{i=1}^{n} g_{i}(x) \leqslant 1\) 对所有 \(x \in X\) 成立，且 \(\sum_{i=1}^{n} g_{i}(x) = 1\) 对 \(x \in K\) 成立。于是对每个函数 \(f \in \mathcal{K}(X, K; \mathbf{C})\) ，可以写 \(f = \sum_{i=1}^{n} fg_{i}\) ，并且有 \(fg_{i} \in \mathcal{K}(X, K_{\alpha_{i}}; \mathbf{C})\) 与 \(\|fg_{i}\| \leqslant \|f\|\) ；若 \(M_{K} = \sum_{i=1}^{n} M_{\alpha_{i}}\) ，则我们得到关系式 (4)。

我们用 \(\mathcal{M}(\mathrm{X};\mathbf{C})\) （如不致混淆，简记作 \(\mathcal{M}(\mathrm{X})\) ）表示 \(\mathrm{X}\) 上的测度所成的向量空间，换言之，即 \(\mathcal{K}(\mathrm{X};\mathbf{C})\) 的对偶。

我们知道，对每个集族 \(\mathfrak{S}\) （由 \(\mathcal{K}(\mathrm{X};\mathbf{C})\) 的有界子集组成的），在 \(\mathcal{M}(\mathrm{X};\mathbf{C})\) 上定义了 \(\mathfrak{S}\) -拓扑，它是局部凸的（TVS, III, §3, No. 1, 命题 1 的推论）。我们把给 \(\mathcal{M}(\mathrm{X};\mathbf{C})\) 赋以 \(\mathfrak{S}\) -拓扑所得到的拓扑向量空间记作 \(\mathcal{M}_{\mathfrak{S}}(\mathrm{X};\mathbf{C})\) 或 \(\mathcal{M}_{\mathfrak{S}}(\mathrm{X})\) 。

命题 7. --- 对 \(\mathcal{K}(\mathrm{X};\mathbf{C})\) 的有界子集所成的、且构成 \(\mathcal{K}(\mathrm{X};\mathbf{C})\) 的覆盖的每个集族 S，空间 \(\mathcal{M}_{\mathfrak{S}}(\mathrm{X};\mathbf{C})\) 是 Hausdorff 的且是拟完备的。

这由 \(\mathcal{K}(\mathrm{X};\mathbf{C})\) 是桶式的这一事实得出（TVS, III, §4, No. 2, Th. 1 的推论 4）。

测度的例子. --- I. 原子测度. 设 X 是一个局部紧空间，a 是 X 的一个点；映射 \(f \mapsto f(a)\) （\(\mathcal{K}(\mathrm{X};\mathbf{C})\) 到 C 中的）显然对 X 的每个包含 a 的紧子集 K 满足条件 (4)（取 \(M_{K} = 1\) ），因而是 X 上的一个测度，记作 \(\varepsilon_{a}\) ；它称为 a 点处的 Dirac 测度，或由放在 a 点处的单位质量所定义的测度。

更一般地，设 \(\alpha\) 是 X 到 C 中的一个映射，使得对 X 的每个紧子集 K，\(\sum_{x\in K}|\alpha(x)|<+\infty\) 。于是，对每个函数 \(f\in\mathcal{K}(X,K;\mathbf{C})\) ，和

\[
\mu (f) = \sum_ {x \in X} \alpha (x) f (x)
\]

有定义，它等于 \(\sum_{x\in K}\alpha (x)f(x)\) ；显然 \(\mu\) 是 \(\mathcal{K}(\mathrm{X};\mathbf{C})\) 上的一个线性形式，且对 \(f\in \mathcal{K}(\mathrm{X},\mathrm{K};\mathbf{C})\) ，

\[
| \mu (f) | \leqslant \left(\sum_ {x \in K} | \alpha (x) |\right) \cdot \| f \|,
\]

换言之，条件 (4) 得到满足。

X 上的测度 \(\mu\) 称为原子测度，如果存在 X 到 C 中的一个映射 \(\alpha\) ，使得对 X 的每个紧子集 K 有 \(\sum_{x\in K}|\alpha (x)| < +\infty\) ，且使得 \(\mu\) 等于如上定义的测度。若 N 是 \(x\in X\) 中使 \(\alpha (x)\neq 0\) 的那些点所成的集，则加于 \(\alpha\) 的条件蕴涵：对 X 的每个紧子集 K，\(\mathrm{K}\cap \mathrm{N}\) 是可数的。也说 \(\mu\) 由质量 \(\alpha (x)\) （放在诸点 \(x\in \mathrm{N}\) 处）所定义。若假定 \(\mathrm{N}\cap \mathrm{K}\) 对每个紧集 \(\mathrm{K}\subset \mathrm{X}\) 是有限的，则显然 \(\sum_{x\in K}|\alpha (x)| < +\infty\) ；这等价于说 N 是 X 的一个闭且离散的子空间，因为这时 X 的每个点有一个只含 N 中有限多个点的紧邻域，反之，若是这种情形，则 X 的每个紧子集可被有限个这样的邻域所覆盖。当 N 是闭且离散的时候，由一个函数 \(\alpha\) （它满足 \(\alpha(x)=0\) ，在 CN 上）所定义的每个原子测度称为 X 上的一个离散测度（参见 §2, No. 5）。

\begin{enumerate}
\def\labelenumi{\Roman{enumi}.}
\setcounter{enumi}{1}
\tightlist
\item
  Lebesgue 测度. 对每个函数 \(f \in \mathcal{K}(\mathbf{R}; \mathbf{C})\) ，存在 R 的一个紧区间 \([a, b]\) ，在其外 f 为零。积分
\end{enumerate}

\[
\operatorname{I} (f) = \int_ {- \infty} ^ {+ \infty} f (x) d x = \int_ {a} ^ {b} f (x) d x
\]

因而有定义；此外，由中值定理（FRV, II, §1, No. 5, 命题 6），有 \(|\mathrm{I}(f)| \leqslant (b - a) \| f \|\) ；这表明 \(f \mapsto \mathrm{I}(f)\) 是 \(\mathbf{R}\) 上的一个测度，称为 Lebesgue 测度。

对 R 的每个区间 J（有界或无界），类似地称测度 \(f \mapsto \int_{J} f(x) dx\) ——\(\mathcal{K}(J; \mathbf{C})\) 上的一个线性形式——为 J 上的 Lebesgue 测度（积分有意义，因为存在含于 J 的紧区间 \([a, b]\) ，在其外 f 为零）。

\begin{enumerate}
\def\labelenumi{\Roman{enumi}.}
\setcounter{enumi}{2}
\tightlist
\item
  设 g 是紧区间 \(I \subset R\) 到 C 中的一个连续映射，在 I 中具有连续导数。设 \(\Gamma = g(I)\) ，它是 C 的一个紧子空间；映射
\end{enumerate}

\[
f \mapsto \int_ {\mathrm{I}} f (g (t)) g ^ {\prime} (t) d t
\]

是 \(\mathcal{C}(\Gamma;\mathbf{C})\) 到 C 中的一个连续线性形式（依中值定理），因而是 \(\Gamma\) 上的一个复测度；关于这一测度的积分也写作 \(\int_{\Gamma} f(z) dz\) ，尽管它不仅依赖于 \(\Gamma\) ，而且依赖于 g。

注. --- 在局部紧空间 X 上给定一个测度 \(\mu\) ，就在 X 上（连同 X 的拓扑一起）定义了一个结构 S。设 \(X_{1}\) 是另一个集合，\(\varphi\) 是 X 到 \(X_{1}\) 上的一个双射；按照一般定义（S, R, §8），结构 \(S_{1}\) ——它由把 X 的结构 S 搬运到 \(X_{1}\) 上（借助 \(\varphi\) ）而得到——按下述方式定义。X 的拓扑被搬运到 \(X_{1}\) 上，所借助的是 \(\varphi\) ；于是 \(\mathcal{K}(X_{1};\mathbf{C})\) 中的函数是使 \(f \circ \varphi\) 属于 \(\mathcal{K}(X;\mathbf{C})\) 的那些函数 f，而测度 \(\mu_{1}\) （\(X_{1}\) 上的）由 \(\mu_{1}(f) = \mu(f \circ \varphi)\) 定义。

特别地，结构 S 的一个自同构是 X 到其自身的一个同胚 \(\sigma\) ，它使得

\[
\mu (f) = \mu (f \circ \sigma)
\]

对每个函数 \(f \in \mathcal{K}(\mathrm{X};\mathbf{C})\) 成立；这时也称测度 \(\mu\) 在同胚 \(\sigma\) 下不变。

例子. --- \(\mathbf{R}\) 上的 Lebesgue 测度在加法群 \(\mathbf{R}\) 的每个平移下不变。事实上，对每个函数 \(f \in \mathcal{K}(\mathbf{R}; \mathbf{C})\) 及每个实数 \(a\) ，有

\[
\int_ {- \infty} ^ {+ \infty} f (x + a) d x = \int_ {- \infty} ^ {+ \infty} f (t) d t
\]

由变量替换公式（FRV, II, §2, No. 1, 公式 (1)）得出。

推广见 Ch. VII, §1, No. 2, Th. 1。

